\DocumentMetadata{
  pdfversion=2.0,
  pdfstandard=ua-2,
  lang=en-US,
  tagging=on,
  tagging-setup={math/setup=mathml-SE,tabsorder=structure}
}
\documentclass[11pt,letterpaper]{article}
\usepackage[margin=0.68in,includefoot]{geometry}
\usepackage{newcomputermodern}
\usepackage{amsmath,amscd,mathtools}
\providecommand{\square}{\mdlgwhtsquare}
\usepackage[hidelinks]{hyperref}
\hypersetup{
  pdftitle={Algebra PhD Exam, May 2016},
  pdfauthor={University of Florida Department of Mathematics},
  pdfsubject={Graduate mathematics examination},
  pdfdisplaydoctitle=true,
  bookmarksopen=true
}
\setcounter{secnumdepth}{0}
\setlength{\parindent}{0pt}
\setlength{\parskip}{0.28em}
\setlength{\emergencystretch}{3em}
\setlength{\leftmargini}{2.15em}
\setlength{\leftmarginii}{2.65em}
\setlength{\leftmarginiii}{3em}
\renewcommand{\labelenumi}{\textbf{\theenumi.}}
\pagestyle{empty}
\raggedbottom
\allowdisplaybreaks
\newenvironment{examproblems}[1]{%
  \begin{enumerate}%
  \setcounter{enumi}{#1}%
  \setlength{\topsep}{0.35em}%
  \setlength{\partopsep}{0pt}%
  \setlength{\itemsep}{0.48em}%
  \setlength{\parsep}{0.18em}%
}{\end{enumerate}}
\newenvironment{examsubparts}{%
  \begin{enumerate}%
  \setlength{\topsep}{0.18em}%
  \setlength{\partopsep}{0pt}%
  \setlength{\itemsep}{0.16em}%
  \setlength{\parsep}{0.1em}%
}{\end{enumerate}}
\newenvironment{examinstructions}{%
  \par\begingroup\itshape
}{%
  \par\endgroup\vspace{0.18em}
}
\makeatletter
\renewcommand\section{\@startsection{section}{1}{0pt}%
  {-1.25ex plus -.25ex minus -.1ex}%
  {0.55ex plus .1ex}%
  {\normalfont\large\bfseries}}
\renewcommand{\@maketitle}{%
  \newpage\null\vskip -1.2em%
  {\footnotesize\bfseries
    \href{https://www.ufl.edu/}{University of Florida}, \href{https://math.ufl.edu/}{Department of Mathematics}\par}%
  \vskip 0.18em%
  \tagstructbegin{tag=Title}%
    {\LARGE\bfseries\href{https://gma.math.ufl.edu/exams/algebra-phd/}{Algebra PhD Exam}, May 2016\par}%
  \tagstructend%
  \vskip 0.35em%
  \tagmcbegin{artifact}%
    \hrule height 1pt%
  \tagmcend%
  \vskip 0.55em%
}
\makeatother
\title{Algebra PhD Exam, May 2016}
\author{}
\date{}
\begin{document}
\maketitle
\thispagestyle{empty}
\begin{examinstructions}
Answer seven problems. Write your answers clearly in complete English sentences. You may quote results (within reason) as long as you state them clearly.
\end{examinstructions}
\begin{examproblems}{0}
\item
Let \(M/K\) be a field extension and let~\(p\) be a prime. Let \(L_1/K, L_2/K\) be finite subextensions of \(M/K\) such that \([L_1:K]\) and \([L_2:K]\) are both powers of~\(p\).
\begin{examsubparts}
\item[\textnormal{(a)}]
Prove that if \(L_1/K\) is a Galois extension then \([L_1L_2:K]\) is a power of~\(p\).
\item[\textnormal{(b)}]
Give an example which shows that if neither of \(L_1/K, L_2/K\) is Galois then \([L_1L_2:K]\) need not be a power of~\(p\).
\end{examsubparts}
\item
Prove that every field~\(K\) is contained in an algebraically closed field.
\item
This problem has two parts.
\begin{examsubparts}
\item[\textnormal{(a)}]
Let~\(C\) be a category and let \(f:X\to Y\) be a morphism in~\(C\). Define what it means for~\(f\) to have each of the following properties:
\begin{examsubparts}
\item[\textnormal{i.}]
\(f\) is a monomorphism.
\item[\textnormal{ii.}]
\(f\) is an epimorphism.
\item[\textnormal{iii.}]
\(f\) is an isomorphism.
\end{examsubparts}
\item[\textnormal{(b)}]
Give an example of a category~\(C\) and a morphism \(f:X\to Y\) in~\(C\) which is a monomorphism and an epimorphism, but not an isomorphism.
\end{examsubparts}
\item
This problem has three parts.
\begin{examsubparts}
\item[\textnormal{(a)}]
Define what it means for a group~\(G\) to be nilpotent.
\item[\textnormal{(b)}]
Prove that if~\(G\) is nilpotent and \(N\trianglelefteq G\) then \(G/N\) is nilpotent.
\item[\textnormal{(c)}]
Prove that if~\(G\) is free on a set~\(S\) and \(|S|>1\) then~\(G\) is not nilpotent.
\end{examsubparts}
\item
Let \(0\to A\mathrel{\mathop{\to}^{f}}B\mathrel{\mathop{\to}^{g}}C\to 0\) be an exact sequence of groups (not necessarily abelian).
\begin{examsubparts}
\item[\textnormal{(a)}]
Suppose there is a group homomorphism \(r:B\to A\) such that \(r\circ f=\operatorname{id}_A\). Prove that \(B\cong A\times C\).
\item[\textnormal{(b)}]
Suppose there is a group homomorphism \(s:C\to B\) such that \(g\circ s=\operatorname{id}_C\). Prove that \(B\cong A\rtimes_\phi C\) for some \(\phi:C\to\operatorname{Aut}(A)\).
\end{examsubparts}
\item
Let~\(R\) be a commutative ring with 1 and let~\(I\) be a proper ideal of~\(R\). The radical of~\(I\) is defined to be 
\[
\sqrt{I}=\{x\in R:x^n\in I\text{ for some }n\geq 1\}.
\]
 Prove that \(\sqrt{I}\) is equal to the intersection of all the prime ideals of~\(R\) which contain~\(I\).
\item
Let~\(R\) be a commutative ring with 1 and let~\(P\) and~\(Q\) be projective~\(R\)-modules. Prove that the~\(R\)-module \(P\otimes_R Q\) is projective.
\item
Let~\(R\) be an integral domain, let~\(K\) be field of fractions of~\(R\), and let \(A\subset K\) be a fractional ideal of~\(R\).
\begin{examsubparts}
\item[\textnormal{(a)}]
Prove that the set \(A'=\{x\in K:xA\subset R\}\) is a fractional ideal of~\(R\).
\item[\textnormal{(b)}]
Prove that the fractional ideal~\(A\) is invertible if and only if \(AA'=R\).
\item[\textnormal{(c)}]
Give an example of an integral domain~\(R\) and a nonzero fractional ideal~\(A\) of~\(R\) which is not invertible.
\end{examsubparts}
\item
Let \(D>1\) be square-free. Determine the integral closure of~\(\mathbb{Z}\) in \(\mathbb{Q}(\sqrt{D})\). Your answer should be carefully justified.
\item
Let~\(R\) be a ring with 1 and let~\(M\) be a unital~\(R\)-module. Assume that for every~\(R\)-submodule~\(N\) of~\(M\) there is an~\(R\)-submodule \(N'\) such that \(M=N\oplus N'\). Prove that if~\(L\) is a nontrivial submodule of~\(M\) then~\(L\) contains a simple submodule.
\item
Determine the number of isomorphism classes of noncommutative semisimple rings with \(2592=2^5\cdot 3^4\) elements. You do not have to list the rings.
\end{examproblems}
\end{document}
